\documentclass[11pt]{article}
\usepackage[margin=1in]{geometry}
\usepackage{amsmath,amssymb,amsfonts,amsthm,mathtools,bm}
\usepackage[T1]{fontenc}
\usepackage{microtype}
\newtheorem{theorem}{Theorem}
\newtheorem{lemma}{Lemma}
\newtheorem{cor}{Corollary}
\newtheorem{defn}{Definition}
\newtheorem{principle}{Principle}
\newtheorem{statement}{Statement}
\theoremstyle{remark}
\newtheorem{statementr}{Statement}

\title{Taylor Expansions of Mappings between Manifolds}
\author{James Taylor}
\date{Working draft derived from the 2003 thesis}
\begin{document}
\maketitle

\section{Taylor expansions of mappings between
manifolds} \label{c:taylor}


\section{Introduction} \label{s:intro}
This chapter is somewhat different from the other chapters. The first
half of the chapter is needed for Chapter \ref{c:algo}. The second
half is answering a natural question that arose while exploring the
first half. 


Let $M$ and $N$ be two smooth manifolds with connections. Let $f: M
\to N$ be smooth. In this chapter, we shall define how to expand $f$
in the best polynomial expansion using the linear structure at hand. 
A connection provides a notion of a local linear
structure on a manifold. Using this linear structure, one can view $f$
as a map between open subsets of vector spaces. Thus, we can expand
$f$ using Taylor polynomials to obtain the Taylor expansion of $f$. 
But a connection also provides a method
for differentiating objects. The question then arises as to whether
using the covariant derivatives in the usual Taylor formula will
provide us with the Taylor expansion of $f$. The answer is no, but it
does when the range is suitably flat.


Given a manifold $M$ and a connection on $M$, $\nabla^M$, we
have a natural identification between the tangent space of a point $p$
and small neighborhoods about $p$ via the exponential map. 
When we wish to view the
neighborhood $U_p \subset M$ as a subset of $T_pM$  using this
identification, then we 
shall refer to  $U_p$ as a \emph{linearized neighborhood};
putting linear coordinates on the tangent space leads to normal
coordinates for our neighborhood. The
identification  allows us to discuss the linearity of
functions. Indeed, let us assume that we have a map $f: M \to N$ and
that $N$ carries a connection, $\nabla^N$. Let $U_p$ be a
linearized neighborhood about $p$ and let $U_{f(p)}$ be a linearized
neighborhood about 
$f(p)$. Then we can discuss linear mappings, quadratic mappings, and,
in general, multilinear mappings between the neighborhoods. It
therefore makes sense to expand $f$ itself into a Taylor series
expansion. The expansion is a sum of multilinear functions, or
\emph{tensors}, evaluated
along the diagonal;  symbolically  we have $f(v) \approx \sum D_T^n
([v]^n).$ We define the \emph{\ensuremath{{n}^{th}}\ Taylor
derivative}\footnote{This is in honor of the usual Taylor expansions
  and not in honor of the author.}
 to be the tensor $D_T^n f$. The \ensuremath{{n}^{th}}\ covariant derivative,
$D_C^n f$, and $D_T^n f $ are both sections of $(T^*M)^n\otimes f_{*}(TN)$,
where $f_* (TN)$ is the pullback of the tangent bundle of $N$ to be
explained later.  
It makes sense to ask whether the symmetric parts agree. This would
imply that using the covariant derivatives in the formulas from
multi-variable calculus would be legitimate. If we have agreement, 
then we say that $f$
satisfies \emph{covariant normality} or is \emph{covariantly normal}.


The following will be proven:  
\begin{enumerate}
\item If $N$ is flat, then all smooth mappings into it are
covariant normality; 
\item If the smooth mapping $f$ maps geodesics to
geodesics, then $f$ is covariantly normal; and
\item Let $N = V$ be a vector bundle
carrying a bundle connection $\nabla^{V}$ over a base manifold $M$ carrying the
manifold connection $\nabla^{M}$. Then those two connections induce a
manifold connection, $\nabla^{N}$, on $N$. Using $\nabla^{N}$, all
smooth sections
of the vector bundle are covariantly normal. 
\end{enumerate}

Results that are about Taylor expansions of sections of a vector bundle
are in the literature, but viewing them as 
Taylor expansion of mappings between
manifolds is not. 


Even
allowing for this weak connection, 
this author could find few papers with a similar flavor to the
material in this chapter.
 An alternative idea to the Taylor expansions for tensors 
is the notion of affine
extensions. Roughly, an affine extension of a tensor 
is defined by choosing normal
coordinates and differentiating the coefficients in that coordinate
system. One then assembles new tensors out of this construction. 
This is different from what we ultimately do in such contexts in the
sense that our basis for the tensor spaces can differ from the basis
that the affine extension concept is based upon.


It was first pursued in the 1920's and 1930's by O. Veblen, T. Thomas,
and H. Ruse. 
More recently, it was discussed in
\cite{GRS98} in the context of symplectic manifolds with
connections. If one is interested in this concept, then that recent
paper  is the paper to learn the material from. It contains references
to the original works as well.  


\section{Covariant derivatives of mappings between manifolds}
\label{s:map}

We start by discussing the standard notions of the pullback bundle
and the pullback 
connection. Given a bundle $E$ over $N$, \ensuremath{{E} \overset{\pi}{\rightarrow} {N}}, and a
manifold $M$ with a smooth map $f: M \to N$, then the pullback
bundle, denoted by \ensuremath{{f}_{*}({E})}, is defined by $\ensuremath{{f}_{*}({E})}_{\ensuremath{q}} \ensuremath{:=}
E_{f(\ensuremath{q})}$. It can be shown that this has the appropriate smooth
structure. Notice that a section of $E$ pulls back to a section of
\ensuremath{{f}_{*}({E})}, but the reverse is only guaranteed if $f$ is a
diffeomorphism. 
If $E$ has a connection, \ensuremath{\nabla^{E}_{} {}}, then the pullback
connection, \ensuremath{\nabla^{fE}_{} {}}, is defined by demanding that the
connection acting on a pulled back section in the direction $X \in TM$ 
should agree with the connection on $E$ acting on the section in the
direction $Df (X)$; $Df : TM \to TN$ is the differential of the map
$f$.  Schematically, the definition demands 
\begin{equation}
\{ \ensuremath{{{\nabla}}_{X}^{fE}{[\ensuremath{{f}_{*}({e})}]}}\}_{p} 
\ensuremath{:=} 
\{\ensuremath{{{\nabla}}_{Df(X)}^{E}{[e]}}\}_{f(p)}
\end{equation}
where $e$ is a section of $E$ and \ensuremath{{f}_{*}({e})}\ is the pull-back
of $e$. 
By choosing a local frame field on $E$ and pulling the frame back, one can
establish what the connection must be in these coordinates.


One
then must show, straightforwardly, that  the different coordinate expressions
describe the same object.
An informative example  is when $f$ maps
$M$ into a single point of $N$. Then the pullback bundle is a trivial
bundle and the pullback connection is a trivial connection.  

Now that we have the general definition of the pullback of a bundle,
we can specialize to the tangent bundles. Connections on the tangent
bundles $TM$ and $TN$ of the manifolds $M$ and $N$ will be denoted by
\ensuremath{\nabla^{M}_{} {}}\ and \ensuremath{\nabla^{N}_{} {}}. Note that if \ensuremath{{M} \overset{f}{\rightarrow} {N}}\ is a
smooth map between $M$ and $N$, then $Df \in \ensuremath{T^* M} \otimes
\ensuremath{{f}_{*}({TN})}$. The Leibniz (product) rule allows us to define the tensor
product of two connections. Thus, we can differentiate $Df$ to obtain
an element of $\ensuremath{T^* M} \otimes \ensuremath{T^* M} \otimes \ensuremath{{f}_{*}({TN})}$. 
Repeating, we can compute the derivatives of all orders of DF.\footnote{ 
It is useful to point out that 
this is 
different then repeatedly taking  differentials. Indeed, the second
differential 
is a map between $T(TM)$ and $T(TN)$ whereas the second covariant
derivative goes from $TM \otimes TM$ to $TN$. 
Using principal bundles, it can be show that there is a relation
between the two versions of derivatives.} This material is from
\cite{EL83}.


The next section discusses the product rule, chain rule, and the
abuses of the tensor world. 
In section \ref{s:key} we discuss the key
observations that makes the entire chapter doable. We then finally discuss
various versions of Taylor expansions of mappings between manifolds. 

\section{The chain rule} \label{s:chain}

The chain rule for derivatives rests on two defining properties of
derivatives: the Leibniz rule and the commutativity of contractions
and derivatives. We first state the main results of this section 
and then explain the
notations and proofs. 


The results of this section are:

\begin{theorem}[Usual Chain Rule]
Let \ensuremath{{M} \overset{F}{\to} {N}
\overset{f}{\to} {P}}\ be smooth maps between connection carrying
manifolds. Let $G$ denote their composition. Then 
\begin{align*}
\ensuremath{D^{n} [{G}] ({v_1, \ldots, v_n}) } = \ensuremath{D^{n} [{G}] ({\mathop{\otimes}_{i=1}^n v_i}) } 
= &\sum_{\ensuremath{Z} \in \ensuremath{\mathcal{\ensuremath{Z}}}(n)} 
\ensuremath{D^{ |\ensuremath{Z}|_R } [{f}] ({\mathop{\otimes}_{r \in R_{\ensuremath{Z}}} \{\ensuremath{D^{|r|} [{F}] ({
\mathop{\otimes}_{c \in C_{r}} v_{\ensuremath{z}_{rc}}}) }\}}) } \\
= & \sum_{k =1}^n \ensuremath{D^{k} [{f}] ({\sum_{\ensuremath{Z} \in \ensuremath{\mathcal{\ensuremath{Z}}}(n,k)}
{\mathop{\otimes}_{r \in R_{\ensuremath{Z}}} \{\ensuremath{D^{|r|} [{F}] ({ 
\mathop{\otimes}_{c \in C_{r}} v_{\ensuremath{z}_{rc}}  }) }\}}}) }.
\end{align*}
\end{theorem}

\begin{theorem}[General Chain Rule]
Let \ensuremath{\mathcal{W}}\ be a vector bundle over $M$, \ensuremath{\mathcal{V}}\ a bundle over $N$,
and \ensuremath{\mathcal{X}}\ a bundle over $P$ and assume that all the bundles have a
connection. Let 
\ensuremath{{M} \overset{F}{\to} {N}
\overset{f}{\to} {P}} be smoot maps between the base manifolds.
Assume $B \in \ensuremath{\Gamma \{ \ensuremath{\mathcal{W}} \otimes \ensuremath{{F}_{*}({\ensuremath{\mathcal{V}}})} \}}$ and $A \in
\ensuremath{\Gamma \{ \ensuremath{{\ensuremath{\mathcal{V}}}^{*}} \otimes \ensuremath{{f}_{*}({\ensuremath{\mathcal{X}}})} \}}$. Then $D[A(B)] = D[A](B,
DF) + A(D[B])$. 
\end{theorem}

Those are the two main theorems. We will use the following two defining
facts. 


\begin{statementr}[Tensor Input]
If A is a multilinear functional acting on $V \times W$, then A may be
viewed as a
linear functional acting on $V \otimes W$. Furthermore, the input
$(v,w)$ into the multilinear A is the same as $v\otimes w$ into the
linear version of A.
\end{statementr}

\begin{statementr}[The Product Rule] 
If we have connections on the bundles $E$ and $F$, both over $M$, then
we define the tensor product connection  on $E
\otimes F$  by demanding the Leibniz rule and commutation
with contractions. Specifically, 
if $A \in \ensuremath{\Gamma \{ E \}}$ and $B \in \ensuremath{\Gamma \{ F \}}$ are any two sections,
then 
\begin{enumerate} 
\item {Leibniz Rule:} 
 $\ensuremath{{{\nabla}}_{X}^{}{[A\otimes B]}} = \ensuremath{{{\nabla}}_{X}^{}{[A]}} \otimes B + A \otimes
\ensuremath{{{\nabla}}_{X}^{}{[B]}}$
\item{Contraction Commutation:} If  $C( A \otimes B)$ indicates any kind of
contraction, then 
$ \ensuremath{{{\nabla}}_{X}^{}{[C(A\otimes B)]}} = C( \ensuremath{{{\nabla}}_{X}^{}{[A \otimes B]}})$. 
\end{enumerate}
\end{statementr}

\begin{proof}[{Diagrams}] {
It is convenient to introduce the notion of a chain rule diagram or 
cr-diagram. It
is a diagram similar to a Young tableaux, but there are important
differences. In a cr-diagram, we left justify the diagrams so that the
first entries lie in a column. The number of boxes in a row has no
particular relation to the other rows i.e.\ the number can 
increase, decrease, or stay the same as we move from row to
row. Furthermore, it is a numbered diagram. We demand that the entries
in the first column strictly increase as we go down the column and the
entries in a row to 
strictly increase as we move to the right. Thus, the number in the
first box in the first column is 1.
The numbers in a column other than the first column
are not constrained to behave in a certain way. In particular, the
numbers in the second column can either increase or decrease. 

Here are several cr-diagrams: {\scriptsize \[ \boxed{\text{12345}} \quad
\boxed{\text{1,2,3}} \quad \boxed{\text{15,2,34}} \quad \boxed{\text{15,24,3}} \quad
\boxed{\text{1,2,345}} \quad \boxed{\text{15,234}}.\]} 
Here are two objects which are not cr-diagrams:
{\scriptsize\[  \boxed{\text{21,345}} \quad \boxed{\text{2,1345}} .\]}


We now define terminologies and notations which will be extremely 
helpful. A cr-diagram will be denoted by \ensuremath{Z}. The length of \ensuremath{Z},
denoted by $ |\ensuremath{Z}| $, is the number of boxes in the diagram.  
The collection of all cr-diagrams will be denoted by
\ensuremath{\mathcal{\ensuremath{Z}}}\ and the set of cr-diagrams of length $n$ will be denoted by
$\ensuremath{\mathcal{\ensuremath{Z}}}(n)$. The 
height of the cr-diagram \ensuremath{Z}\ is the number of rows and it will be
denoted by $ |\ensuremath{Z}|_{R} $. The set of cr-diagrams of length $n$ and height
$m$, shall be denoted by $\ensuremath{\mathcal{\ensuremath{Z}}}(n,m)$. 
The entry in row $r$ and column $c$ will be
denoted by $\ensuremath{z}_{rc}$. The length of  row $r$ will be denoted by
$ |r| $. The set of rows in $\ensuremath{Z}$ will be denoted by $R_{\ensuremath{Z}}$. The
set of columns in row  $r$ will be $C_{ r}$. The diagram 
{\scriptsize $ \ensuremath{Z}= {\boxed{\text{147,2356}}}$} 
has the
following properties: $ |\ensuremath{Z}| =7 $, $ |\ensuremath{Z}|_R = 2$, $ \ensuremath{z}_{12} = 4$.  
Note that the cr-diagram of length $n$ and height $n$ is unique and is
a single column. Also unique is the single row diagram which is of
length $n$ and height $1$.


A product rule diagram (pr-diagram) is a much simpler diagram. Each row must
be strictly increasing, but there are no constraints for the first
column. We also need to allow for rows with no entries. We do that by
putting one box for the row and putting a zero in that entry. We use
the same notations as before for the cr-diagram except we shall denote 
such objects by $PR$ and the set of product rule diagrams with $m$
rows as $\ensuremath{\mathcal{PR}}_m$. Notice that the number of rows, which includes
empty rows, does not change as one increases the number of derivatives.
}
 \end{proof}
\begin{proof}[{Notations}] {
We shall use $[ \hspace{.05in}]$ to denote an object which is being
differentiated. Hence, functions should not freely pass
through. Parentheses shall be mainly used to indicate some kind of
 contraction. The
simplest case is  $A(B) = B(A)$ when $A \in \ensuremath{{E}^{*}}$ and $ B \in E$, then
this is just usual composition. This also defines the kernel of the
contraction 
operation. Namely, \[C(A \otimes B) \ensuremath{:=} A(B).\] As contractions can
be more general and ambiguous, caution and excellent notation is
required. A discussion of what we shall eventually need is in the
appendices. For other uses of parenthetical expressions, we shall use
$\{ \}$. 


We shall often employ
$D^n[A]$ to denote the \ensuremath{{n}^{th}}\ covariant derivative acting on $A \in E$,
where $E$ is some vector bundle. Note that to define the higher
covariant derivatives, we essentially need a connection on the base
space. Namely, 
$\ensuremath{D^{} [{A}] {} } \in E \otimes \ensuremath{T^* M}$ and therefore $D^n[A]$ is 
defined via the Leibniz rule 
if we have a connection on the bundle and its base space. Note
that $\ensuremath{D^{} [{A}] ({X}) } \ensuremath{:=} \ensuremath{{{\nabla}}_{X}^{E}{[A]}}$ and $\ensuremath{D^{} [{A}] ({B \otimes X}) }
\ensuremath{:=} B(\ensuremath{{{\nabla}}_{X}^{E}{[A]}}) = \ensuremath{{{\nabla}}_{X}^{E}{[A]}}( B)$. Furthermore,
$D_p [A]$ will denote
 the covariant derivative evaluated at $p \in M$. It is
useful to remark that $\ensuremath{D^{2} [{A}] ({X\otimes Y}) } \neq
\ensuremath{{{\nabla}}_{Y}^{}{[\ensuremath{{{\nabla}}_{X}^{}{[A]}}]}}$ as the left-hand side is a tensor in
both $X$ and $Y$ whereas the right-hand side is only a tensor in
$Y$.

The set of sections of a bundle $E$ shall be denoted by
\ensuremath{\Gamma \{ E \}}. Recall that a section of a bundle is a smooth map from
the base space into the bundle such that the projection map acting on
the section is the identity map on the base space.  

The tensor product will occur often and the ordering of the elements
will be crucial. Therefore, unless otherwise noted, when we write the
tensor product over an ordered set, we shall mean the tensor product
in that order. In other words, 
\[
\mathop{\otimes}_{i = 1}^n v_i \ensuremath{:=} v_1 \otimes v_2 \otimes \dotsm \otimes v_n
\]
}
 \end{proof}


\begin{proof}[{Comment on tensor input}] {


The fact to note is that the tensor takes over the role of
multilinearity, enabling the function A to be viewed as linear. Notice
that once we view the input as a tensor product, we can use the
Leibniz rule on it. }
 \end{proof}


\begin{proof}[{Comment on product rule}] {
 The product rule allows us to define the 
covariant derivative on the dual bundle since, for $A \in
 \ensuremath{\Gamma \{ \ensuremath{{E}^{*}} \}}$ and $B \in \ensuremath{\Gamma \{ E \}}$, they must satisfy 
$\ensuremath{D^{} [{A}] ({B}) } = \ensuremath{D^{} [{A(B)}] {} } - A(\ensuremath{D^{} [{B}] {} })$. As $A(B)$ should be
 a function over $M$, we know its derivative and, by assumption, the
 derivative is defined on $E$. In particular a connection on the
 tangent bundle provides a connection on all of the tensor product
 bundles involving only the tangent and cotangent bundles of $M$.


}
 \end{proof} 


\begin{proof}[Proof of the general chain rule]
Since $A$ is a section of a bundle over $N$, we can pull it back to
the pull-back bundle over $M$. Thus, \[\ensuremath{{F}_{*}({A})} \in
\ensuremath{{F}_{*}({\ensuremath{\Gamma \{ \ensuremath{{\ensuremath{\mathcal{V}}}^{*}} \otimes \ensuremath{{f}_{*}({\ensuremath{\mathcal{X}}})} \}}})} \subset
\ensuremath{\Gamma \{ \ensuremath{{\ensuremath{{F}_{*}({\ensuremath{\mathcal{V}}})}}^{*}} \otimes \ensuremath{{(fF)}_{*}({\ensuremath{\mathcal{X}}})} \}}. \]  
Then $A(B)$ may be understood as contracting over the pulled back
 $\ensuremath{\mathcal{V}}$ bundles. In
particular, we have that $\ensuremath{{{\nabla}}_{X}^{M}{[A(B)]}} = \ensuremath{{{\nabla}}_{X}^{M}{[\ensuremath{{F}_{*}({A})}]}}(B) +
A(\ensuremath{{{\nabla}}_{X}^{M}{[B]}})$. This is just the product rule for tensor product
bundles. Since $\ensuremath{{{\nabla}}_{X}^{M}{[\ensuremath{{F}_{*}({A})}]}} \ensuremath{:=}
\ensuremath{{{\nabla}}_{\ensuremath{D^{} [{F}] ({X}) }}^{N}{[A]}}$, we actually have our result.
\end{proof}


\begin{proof}[Proof of the usual chain rule]
The starting
case is the usual well-known chain rule.\footnote{The method of proof
depends on the definition of tangent vector used. Coordinates and the
multi-variable chain rule can always be used. If we use the notion of
tangent vectors as equivalence classes of curves, then $D[F]$ is defined
by taking $ v= [\gamma] \to [F(\gamma)] \ensuremath{:=} D[F](v)$; it is necessary to
check that 
this is independent of the representative $\gamma$. But then the basic
chain rule is simply composition: $ [\gamma] \to [G(\gamma))] = D[G](v)$ is the
same as $[ \gamma] \to [F \gamma] = D[F](v)  \to [f(F(\gamma))]=  
D[f](D[F](v))$.} 
Our inductive step goes as follows. We assume  
\begin{equation*}
\ensuremath{D^{n} [{G}] ({\mathop{\otimes}_{i=1}^n v_i}) } 
= \sum_{\ensuremath{Z} \in \ensuremath{\mathcal{\ensuremath{Z}}}(n)} 
\ensuremath{D^{ |\ensuremath{Z}|_R } [{f}] ({ \mathop{\otimes}_{r \in R_{\ensuremath{Z}}} \{\ensuremath{D^{ |r| } [{F}] {} }
(\mathop{\otimes}_{c \in C_{r}} v_{\ensuremath{z}_{rc}} )\}}) } 
\end{equation*}


Let us take the $v_i$ to be vector fields and differentiate the
resultant scalar field. The terms involving the derivatives of the
vector field will cancel while the other remaining terms will be what
we are interested in.
The calculations follows from the product rule and general
chain rule. 
We start with what we want to find which is

\begin{align*}
\ensuremath{{{\nabla}}_{v_{n+1}}^{}{[\ensuremath{D^{n} [{G}] ({\mathop{\otimes}_{i=1}^n v_i}) }]}} &=
\ensuremath{{{\nabla}}_{v_{n+1}}^{}{[\ensuremath{D^{n} [{G}] {} }]}}(\mathop{\otimes}_{i=1}^n v_i) 
+ \ensuremath{D^{n} [{G}] ({\ensuremath{{{\nabla}}_{v_{n+1}}^{}{[\mathop{\otimes}_{i=1}^n v_i]}}}) } \\
 &=
{\ensuremath{D^{n+1} [{G}] {} }}(\mathop{\otimes}_{i=1}^{n+1} v_i) 
+ \ensuremath{D^{n} [{G}] ({\ensuremath{{{\nabla}}_{v_{n+1}}^{}{[\mathop{\otimes}_{i=1}^n v_i]}}}) };
\end{align*}

the second term is an extra complication which will also appear
in the following:

\begin{align*}
\ensuremath{{{\nabla}}_{v_{n+1}}^{}{[\ensuremath{D^{n} [{G}] ({\mathop{\otimes}_{i=1}^n v_i}) }]}}
& = \sum_{\ensuremath{Z} \in \ensuremath{\mathcal{\ensuremath{Z}}}(n)} 
\ensuremath{{{\nabla}}_{v_{n+1}}^{}{[\ensuremath{D^{ |\ensuremath{Z}|_R } [{f}] ({ \mathop{\otimes}_{r \in R_{\ensuremath{Z}}}
\{\ensuremath{D^{|r|} [{F}] {} } 
( \mathop{\otimes}_{c \in C_{r}} v_{\ensuremath{z}_{rc}} )\}}) } ]}} \\
&=
\sum_{\ensuremath{Z} \in \ensuremath{\mathcal{\ensuremath{Z}}}(n)} 
\ensuremath{{{\nabla}}_{v_{n+1}}^{}{[\ensuremath{D^{ |\ensuremath{Z}|_R } [{f}] {]} }}}
{( \mathop{\otimes}_{r \in R_{\ensuremath{Z}}} \{\ensuremath{D^{ |r| } [{F}] {} }
( \mathop{\otimes}_{c \in C_{r}} v_{\ensuremath{z}_{rc}} )\})} \\
&\quad+
\sum_{\ensuremath{Z} \in \ensuremath{\mathcal{\ensuremath{Z}}}(n)} 
\ensuremath{D^{ |\ensuremath{Z}|_R } [{f}] {(} }\sum_{\tilde r \in R_{\ensuremath{Z}}} 
 \{\mathop{\otimes}_{r \in R_{\ensuremath{Z}}, r\neq \tilde r} 
\ensuremath{D^{|r|} [{F}] {} }
(\mathop{\otimes}_{c \in C_{r}} v_{\ensuremath{z}_{rc}} ) \} 
\\& \hspace{1.5in}\otimes
\ensuremath{{{\nabla}}_{v_{n+1}}^{}{[  
\ensuremath{D^{|\tilde r|} [{F}] {} }]}}
(\mathop{\otimes}_{c \in C_{\tilde r}} v_{\ensuremath{z}_{\tilde rc}} ) ) \\
&\quad+
\sum_{\ensuremath{Z} \in \ensuremath{\mathcal{\ensuremath{Z}}}(n)} 
\ensuremath{D^{ |\ensuremath{Z}|_R } [{f}] {(} }\sum_{\tilde r \in R_{\ensuremath{Z}}} 
 \{\mathop{\otimes}_{r \in R_{\ensuremath{Z}}, r\neq \tilde r} 
\ensuremath{D^{|r|} [{F}] {} }
(\mathop{\otimes}_{c \in C_{r}} v_{\ensuremath{z}_{rc}} ) \} 
\\& \hspace{1.5in}\otimes  
\ensuremath{D^{|\tilde r|} [{F}] {} }
(\mathop{\otimes}_{c \in C_{\tilde r}} \ensuremath{{{\nabla}}_{v_{n+1}}^{}{[
  v_{\ensuremath{z}_{\tilde{r}c}} ]}}) ) \\ 
&\hspace{-.45in}=
\sum_{\ensuremath{Z} \in \ensuremath{\mathcal{\ensuremath{Z}}}(n)} 
\ensuremath{D^{ |\ensuremath{Z}|_R +1} [{f}] {(} }\{ \mathop{\otimes}_{r \in R_{\ensuremath{Z}}} \ensuremath{D^{ |r| } [{F}] {} }
( \mathop{\otimes}_{c \in C_{r}} v_{\ensuremath{z}_{rc}} )\} \otimes \ensuremath{D^{} [{F}] ({v_{n+1}}) }) \\
&\quad+
\sum_{\ensuremath{Z} \in \ensuremath{\mathcal{\ensuremath{Z}}}(n)} 
\ensuremath{D^{ |\ensuremath{Z}|_R } [{f}] {(} }\sum_{\tilde r \in R_{\ensuremath{Z}}} 
 \{\mathop{\otimes}_{r \in R_{\ensuremath{Z}}, r\neq \tilde r} 
\ensuremath{D^{|r|} [{F}] {} }
(\mathop{\otimes}_{c \in C_{r}} v_{\ensuremath{z}_{rc}} ) \} 
\\& \hspace{1.25in}\otimes
\ensuremath{{{\nabla}}_{v_{n+1}}^{}{[  
\ensuremath{D^{|\tilde r|} [{F}] {} }]}}
(\{\mathop{\otimes}_{c \in C_{\tilde r}} 
v_{\ensuremath{z}_{\tilde rc}}\} \otimes v_{n+1}) ) \\
&\quad+
\ensuremath{D^{n} [{G}] ({\ensuremath{{{\nabla}}_{v_{n+1}}^{}{[\mathop{\otimes}_{i=1}^n v_i]}}}) }.
\end{align*}

The first term, in the set of three terms, requires the use of the
general chain rule.  That set of terms 
 corresponds to cr-diagrams of length $n+1$
obtained from those of length $n$ by attaching a new row at the
bottom; the new row 
contains exactly 
one box whose entry is $n+1$. The second set of terms requires the use
of the product
rule. It corresponds to modifying the old cr-diagrams by appending a
new box to the end of a row. For every cr-diagram of length $n$ and
height $l$, we obtain $l$ new cr-diagrams with the same height but
with length $n+1$ and $n+1$ is again the entry in the new box. The
third set of terms is an extraneous set for which the product rule and
linearity is of crucial importance for the manipulation of those
terms. 
\end{proof}


\section{The primary tools} \label{s:key}

In this section, we prove that a certain group of 1-forms, associated
to a point p, have the property that the fully symmetrized version of the
covariant derivatives of those forms are 
zero at the point $p$. These forms  are the
coordinate forms in normal coordinates. We then use them to prove that
the fully symmetrized covariant derivatives of any $(0,l)$ tensor is
given by Euclidean differentiation in normal coordinates.

\begin{theorem}\label{tt:key}
Let \ensuremath{\gamma_v}\ be a geodesic. Assume that both $\ensuremath{{{\nabla}}_{v}^{}{[A(B)]}}=
0$ and $\ensuremath{{{\nabla}}_{v}^{}{[B]}} = 0$ along
\ensuremath{\gamma_v}. Then \ensuremath{D^{n} [{A}] ({B, v^n}) } = 0  along \ensuremath{\gamma_v}.    
\end{theorem}

\begin{cor}\label{tc:da}
Let \ensuremath{\gamma_v}\ be a geodesic. Assume that $\ensuremath{{{\nabla}}_{v}^{}{[A]}}=
0$ along
\ensuremath{\gamma_v}. Then \ensuremath{D^{n} [{A}] ({ v^n}) } = 0  along \ensuremath{\gamma_v}.   
\end{cor}

\begin{cor}\label{tc:geo}
Let $(M,\nabla^{M})$ be a smooth manifold with connection. Let 
$(x_1,
\dots, x_m)$ be the coordinate functions for 
 normal coordinates centered at $p \in M$. Then
\ensuremath{D_{p}^{n} {[dx^i]} ({v^{n+1}}) } = 0.
\end{cor}

\begin{cor}\label{tc:imp}
Let $(M,\nabla^M)$ be a smooth manifold with connection.  Let 
$(x_1,
\dots, x_m)$ be the coordinate functions for 
 normal coordinates centered at $p \in M$. 


Let
$\omega$ be a $(0,l)$ tensor. If we view the coordinates as providing a
linear structure, we can compute the $\ensuremath{{n}^{th}}$ 
Euclidean derivative of $\omega$, $\tilde{D}^n_{p}[\omega]$.
Then
$\ensuremath{D_{p}^{n} {[\omega]} ({v^{l+n+1}}) } = \tilde{D}^n_{p}[\omega](v^{l+n+1})$.
\end{cor}

\begin{proof}[Proof of Theorem~\ref{tt:key}]
It is a simple induction using the product rule. Assume that
$\ensuremath{D^{n-1} [{A}] ({B, v^{n-1}}) } = 0$ along \ensuremath{\gamma_v}. Notice $n=1$
is our assumption. Then 


\begin{align*}
(\ensuremath{{{\nabla}}_{v}^{}{[\ensuremath{D^{n-1} [{A}] {} }]}})(B, v^{n-1})  =&
\ensuremath{{{\nabla}}_{v}^{}{[(\ensuremath{D^{n-1} [{A}] {(B, v^{n-1})} })]}} \\ 
&- \ensuremath{D^{n-1} [{A}] ({\ensuremath{{{\nabla}}_{v}^{}{[B]}}, v^{n-1}}) }
- \ensuremath{D^{n-1} [{A}] {(B, \ensuremath{{{\nabla}}_{v}^{}{[v^{n-1}]}})} }.
\end{align*}


By assumption, the objects being differentiated along the geodesic is 
zero along the geodesic; thus, the first two terms vanish. 
first and second terms are zero. Since, along a geodesic,
$\ensuremath{{{\nabla}}_{v}^{}{[v]}} = 0$ and $v^{n-1}$ is the $(n-1)$-fold tensor power of
$v$ with itself, 
we have that the third term is zero by the Leibniz rule. It is the
vanishing of the third term which requires $\ensuremath{\gamma_v}$ to be a geodesic.
\end{proof}

\begin{proof}[Proof of Corollary~\ref{tc:da}]
Since $A(\ensuremath{\text{Id}}) = A$, we may apply Theorem~\ref{tt:key} with $B=\text{Id}$.  
All we need to establish
is that $\ensuremath{D^{} [{\text{Id}}] {} } =0 $. The quick statement is that this
must be true for consistency's sake. Indeed, 
\[\ensuremath{D^{} [{\ensuremath{\text{Id}}(\ensuremath{\text{Id}})}] {} } = \ensuremath{D^{} [{\ensuremath{\text{Id}}}] ({\ensuremath{\text{Id}}}) } + \ensuremath{\text{Id}}(\ensuremath{D^{} [{\ensuremath{\text{Id}}}] {} }.\]
Since it is the identity, this says that 
\[\ensuremath{D^{} [{\ensuremath{\text{Id}}}] {} } = 2 \ensuremath{D^{} [{\ensuremath{\text{Id}}}] {} }\]
which immediately implies that it is a zero 
tensor.\footnote{A coordinate proof would go as follows. Pick a frame
  field, say 
 $\{e_i\}$,  and its dual frame field $\{e^i\}$. Then we
have that $e^i(e_j)$ is a constant function in the neighborhood and
therefore its 
derivative is zero. Thus, the product rule implies that
$\ensuremath{{{\nabla}}_{v}^{}{[e^i]}}(e_j) = -e^i(\ensuremath{{{\nabla}}_{v}^{}{[e_j]}})$. Since 
 $I = e^i \otimes e_i$, we have 
\begin{align*}
\ensuremath{D^{} [{\ensuremath{\text{Id}}}] ({e_j, v}) } & = \\
& = \sum_i\{ (\ensuremath{D^{} [{e^i}] ({e_j, v}) }) \otimes e_i + 
(e^i)(e_j) \otimes (\ensuremath{D^{} [{e_i}] ({v}) })\} \\
& = \sum_i\{-e^i(\ensuremath{D^{} [{e_j}] ({v}) }) \otimes {e_i}\} +
 (\ensuremath{D^{} [{e_j}] ({v}) }) \\
& = - \ensuremath{\text{Id}}(\ensuremath{D^{} [{e_j}] ({v}) }) + \ensuremath{D^{} [{e_j}] ({v}) } = 0.
\end{align*}
Thus we have that the identity tensor always has a zero covariant
derivative.}
\end{proof}


\begin{proof}[Proof of Corollary~\ref{tc:geo}]
In normal coordinates, the image of $v \in T_p M$ under the
exponential map $\ensuremath{\exp_{p}{}} : T_p M \to M$ is labelled $v$. 
We also have 
$tv = \ensuremath{\exp_{p}{tv}}= \gamma_{tv}(1) = \gamma_{v}(t)$,  
( $t$ between 0 and 1). 
In particular, \ensuremath{{\overset{\cdot}{\gamma}}_{v}({t})} is simply $v$ where \ensuremath{\gamma_v}\ is a geodesic
starting at $p$ with $\ensuremath{{\overset{\cdot}{\gamma}}_{v}({0})} = v$.  
As \ensuremath{{\overset{\cdot}{\gamma}}_{v}({t})} is constant along \ensuremath{\gamma_v}, it can be extended to a
constant vector field in
the normal neighborhood. Thus, $dx^i(v)$ is a constant
function. We have $\ensuremath{{{\nabla}}_{v}^{}{[dx^i(v)]}} = 0 $ and $\ensuremath{{{\nabla}}_{v}^{}{[v]}}
=0$. Since $\ensuremath{{{\nabla}}_{v}^{}{[dx^i]}} (v) =  
\ensuremath{{{\nabla}}_{v}^{}{[(dx^i(v))]}} + dx^i(\ensuremath{{{\nabla}}_{v}^{}{[v]}})=0$, we are in the
situation where Theorem~\ref{tt:key} applies. Thus, we have that
$\ensuremath{D^{n} [{dx^i}] ({v, v^n}) } = 0$ along \ensuremath{\gamma_v}\ and, in particular, at $p$. 
\end{proof}

\begin{proof}[{Warning}] { Although the parallel transport of $v$ along its geodesic is
simply a constant in these coordinates, in general the other
coordinate vector fields will
not have this property on \ensuremath{\gamma_v}. If they did, then this would be a
locally trivial
connection. 
Indeed, if the connection was the Levi-Civita connection for a metric,
 then, as the metric is preserved by
parallel transport, the metric would be a flat metric. 
}
 \end{proof}

\begin{proof}[Proof of Corollary~\ref{tc:imp}]
Using coordinates, we have  $\omega \ensuremath{:=} \omega_{i_1 \cdots i_n}
dx^{i_1} \cdots dx^{i_n}$ where summation is implied over repeated
indices. 
To compute the derivatives, we use the product rule repeatedly. 
Indeed, we can take the above coordinate expression and rewrite it as 
$\omega_{i_1 \cdots i_n}\otimes
dx^{i_1} \otimes\cdots \otimes dx^{i_n}$. From Corollary~\ref{tc:geo},
we have that when we evaluate $\ensuremath{D^{n} [{\omega}] ({v^{n+l+1}}) }$, the terms
involving the derivatives of the coordinate forms vanish. Thus, we are
only left with the derivative of the coefficient, which is exactly
what the Euclidean derivative is. 
\end{proof}
 


\section[Vector expansions]{Taylor expansions of mappings between
finite dimensional  
vector spaces} \label{s:tay}

In this section, we shall formulate the question properly. We start by
recalling the usual facts about the Taylor expansion in vector
spaces which may be found, for example, in \cite{AMR89}. 
The \ensuremath{{n}^{th}}\ Taylor function of $f$, $T^n f$, is the best $n$-polynomial
approximation to 
$f$ up to terms of order $o(\epsilon^n)$. 

This means, precisely, the following. We start with the finite
dimensional vector spaces
$V$, $W$, and a smooth map $f:V \to W$. Then the Taylor
function $T^nf$ at $x_0$ 
should be an \ensuremath{{n}^{th}}-degree polynomial,
i.e. 
\[
T^nf \in \{\bigoplus_{i=0}^n (\ensuremath{{V}^{*}})\ensuremath{^{\otimes i} }\} \otimes W.\footnote{
A vector space raised to the \ensuremath{{0}^{th}}-power is defined to be the field
itself. Also, polynomials are always taken to be symmetric
objects. That is to say, one mainly wants to evaluate them on the diagonal.}
\]
We require that $T^nf$ satisfies 
\begin{equation} \label{te:tay}
 \lim_{x \to x_0} 
\frac{\lVert f(x) - T^nf(\oplus_{i=0}^n \{x-x_0\}^i) \rVert_W } 
{\lVert x-x_0 \rVert_V^{n}} = 0
\end{equation}


This statement is
independent of which norm is chosen since, in finite dimensions, all
norms are equivalent.

We shall take $x_0$ to be 0 and
write $v$ for the vector $x-x_0$. 


We first establish that the symmetric part is unique. Indeed,
let $A$ and $B$ be two polynomials satisfying \eqref{te:tay}. We will be
done if we show that for any vector $v \in V$, 
\[(A-B)(\oplus_{i=1}^n \{v\}^i) = 
(A-B)(v \oplus \{v\otimes v\} \oplus \dotsb \oplus \{\otimes_{i=1}^n v
 \})
 = 0. \] 
 But this is easy. We first note that 
\begin{equation} \label{te:unt}
\lim_{v \to 0} 
\frac{\lVert (A-B)(v) \rVert_W } 
{\lVert v \rVert_V^{n}} = 0;
\end{equation}
this follows by inserting 0 in the form of $f-f$ into the numerator, 
using the triangle inequality, and then using \eqref{te:tay}.
Let $A_i$ and $B_i$ be the $i$-linear parts of $A$ and $B$,
respectively. 


By replacing $v$ with
$hv$, $h \in \mathbf{R}$, we have 
\begin{equation} 
0 = \lim_{h \to 0} 
\frac{\lVert \sum_{i=1}^n (A_i-B_i)(\{hv\}^i) \rVert_W } 
{\lVert hv \rVert_V^{n}}  = \lim_{h \to 0^+} 
\frac{ \lVert \sum_{i=1}^n \frac{1}{h^{n-i}} (A_i-B_i)(v^i) \rVert_W } 
{ \lVert v \rVert_V^{n}}.
\end{equation}
Thus, $\lVert (A_i-B_i)(v^i) \rVert_W$ must be zero for every $i$
which implies  uniqueness. 

We now prove that 
\[
T^n f (v) = \sum_{i=0}^n \frac{1}{i!} \ensuremath{D_{0}^{i} {[f]} ({v}) }.
\] 
All we need
is the fundamental theorem of calculus and integration by parts. The
fundamental theorem of calculus is   
\[ f(u + v) - f(u) = \int_{0}^1 \ensuremath{D_{u+tv}^{} {[f]} ({v}) } dt.\]
We now use integration by parts which is the reverse of the product
rule. 
Specifically, we claim that 
\[
\int_{0}^1 \frac{(1-t)^{i-1}}{(i-1)!} \ensuremath{D_{u+tv}^{i} {[f]} ({v^{i}}) } dt = 
\frac{1}{i!} \ensuremath{D_{u}^{i} {[f]} ({v^{i}}) } + 
\int_{0}^1 \frac{(1-t)^i}{i!} \ensuremath{D_{u+tv}^{i+1} {[f]} ({v^{i+1}}) }   dt.
\]

This follows immediately from 
\[\begin{split}
0 - (-\frac{1}{i!} \ensuremath{D_{u}^{i} {[f]} ({v^{i}}) }) &= \int_{0}^1
\ensuremath{\frac{d{}}{d{t}}}[\frac{-(1-t)^i}{i!}
\ensuremath{D_{u+tv}^{i} {[f]} ({v^{i}}) }] dt \\
&=  \int_{0}^1 \frac{(1-t)^{i-1}}{(i-1)!} \ensuremath{D_{u+tv}^{i} {[f]} ({v^{i}}) } dt 
\\&\quad+ \int_{0}^1 \frac{-(1-t)^i}{i!} \ensuremath{D_{u+tv}^{i+1} {[f]} ({v^{i+1}}) }   dt.
\end{split}
\]
Since $f$'s  derivatives are continuous both as operator-valued
functions over $M$, we have that 
\[
\lVert D^{l}_{q} [f] (v^{l}) \rVert_W \leq C_l \lVert v^{l}  \rVert_V
\]
for  $q$ in a small enough neighborhood.
This immediately leads to  
\[
{\lVert f(x) - \sum_{i=0}^n \frac{1}{i!} \ensuremath{D_{0}^{i} {[f]} ({v}) } \rVert_W }
=  \lVert \int_{0}^1 \frac{(1-t)^n}{n!} \ensuremath{D_{u+tv}^{n+1} {[f]} ({v^{n+1}}) }
dt
\rVert_W \leq K_{n+1} \lVert v \rVert_V^{n+1}  
\]
for some constant $K_{n+1}$ and $v$ in the appropriate neighborhood of
the origin.
Thus, we have our limit. 

Having discussed the vector space case, we now formulate
the question for manifolds. If $M$ and $N$ have connections, then
there is a canonical identification via the exponential map 
between a neighborhood of $0\in
T_pM$ and a neighborhood of $p \in M$. The same is true for $N$. Thus
any mapping $f$ between $M$ and $N$ can be viewed as a mapping between
some neighborhood of $0 \in T_pM$ and $T_{f(p)}N$. We can then Taylor
expand that map. Since these are vector spaces, their tangent spaces
are canonically identifiable with themselves. Thus, we end up with a
polynomial mapping from the tangent space of $M$ to that of the
tangent space of $N$. 
The Taylor polynomial of $f$ is the polynomial between the tangent
spaces obtained through the Taylor expansion of $f$ as a mapping
between the two tangent spaces. The \ensuremath{{i}^{th}}\ monomial of the
polynomial is indepedent of $n$; we shall call this the \ensuremath{{i}^{th}}\
Taylor derivative.
We could also construct a polynomial out of the
covariant derivatives of these mappings. Are these polynomials the
same? That is the question. The answer is that, loosely
speaking, they  agree when the range is appropriately 
flat. 

A map $f$ is \emph{covariantly normal at $p$} iff its
symmetrized covariant derivatives at $p$ agree with the Taylor
derivatives. If $f$ is covariantly normal at all points of $M$, then
$f$ is \emph{covariantly normal}. 

\section{Taylor expansions of mappings into flat spaces}

In this section we discuss mappings into a flat space. As we shall
prove, these maps are always covariantly normal. 


Let $(N, \nabla^N)$ 
be a flat cc-manifold. This means that both the torsion and the
curvature of the connection vanishes. In this section, the following
standard lemma is of crucial importance.

\begin{lemma} \label{tl:flat}
Let $(N, \nabla^N)$ 
be a flat cc-manifold and $(U, (x^1, \ldots x^n))$ be normal
coordinates for $U$ centered at $p$. 
Then $\nabla^N [\ensuremath{\frac{\partial {}}{\partial {x^i}}}] =0 $ everywhere in
$U$ for every $i$, i.e.\ the coordinate vector fields are parallel.  
\end{lemma}


\begin{proof}
Let $(U, (y_1,\ldots, y_n))$ be a coordinate neighborhood such that 
the coordinate vector fields are parallel and such that $\ensuremath{\frac{\partial {}}{\partial {y^i}}}\vert_p =
\ensuremath{\frac{\partial {}}{\partial {x^i}}}\vert_p$. Existence of such a
neighborhood follows from 
flatness as in the proof of Theorem 7.3 in the text \cite{Lee98}. 
The claim is that $y^i = x^i$. As the $x^i$'s are normal coordinates,
we know, by definition, that in these coordinates 
 $\gamma (t) \ensuremath{:=} t\sum_{i=1}^n a_i x^i$
is a geodesic such that $\gamma(0) = 0 (=p)$ and 
 $\dot{\gamma}(0) = a_i \ensuremath{\frac{\partial {}}{\partial {x^i}}}\vert_p$. But if we define
$\phi(t) \ensuremath{:=} t\sum_{i=1}^n a_i y^i$, then we claim that this is a
geodesic.  Once we establish that claim, then we are done by
uniqueness results since
$\gamma$ and $\phi$ do have the same initial conditions; hence, the
coordinates are the same (the point represented by $\sum a_i x^i$ is
the same, through the geodesics, as $\sum a_i y^i$).

We now need to establish the $\phi(t)$ is a geodesic. We thus need to
show that $\ensuremath{{{\nabla}}_{\dot{\phi}}^{N}{[\dot{\phi}]}} = 0$. We know that
$\dot{\phi} (t) = \sum_{i=1}^n a_i \ensuremath{\frac{\partial {}}{\partial {y^i}}} \vert_{\phi(t)}$. Thus
the vector field along the curve agrees with the vector field
$\sum_{i=1}^n a_i \ensuremath{\frac{\partial {}}{\partial {y^i}}}$. And therefore, by definition, 
$\ensuremath{{{\nabla}}_{\dot{\phi}}^{N}{[\dot{\phi}]}} = \ensuremath{{{\nabla}}_{\sum_{i=1}^n a_i
\ensuremath{\frac{\partial {}}{\partial {y^i}}}}^{N}{[\sum_{i=1}^n a_i \ensuremath{\frac{\partial {}}{\partial {y^i}}}]}}$. But this is zero as all
of the $\ensuremath{\frac{\partial {}}{\partial {y^i}}}$ are parallel and the $a_i$ are constants. 
\end{proof}

The above proof was rather indirect. Indeed, in the proofs of the
existence of the $y$ coordinates, all of them that the author  has seen 
requires the construction
of a new coordinate system out of a given coordinate
system. Even more galling,  in \cite{Lee98}, one starts with some
coordinate system; that author suggests that normal coordinates is a
possible choice. Then the proof constructs that coordinate system
using theorems about integrability of systems. But, as we just showed,
this constructs normal coordinates. Thus, one uses an elaborate
machinery to construct coordinates that one already has. 

With the lemma, we may now prove the following theorem.
\begin{theorem} \label{t:sim}  
Let f be a map from $(M, \nabla^M)$ into the flat cc-manifold
$(N, \nabla^N)$. 
Then $f$ is covariantly normal. 
\end{theorem}

\begin{proof}
Choose normal coordinate neighborhoods
about $p \in M$ and $f(p) \in N$. 
Using these coordinates, we write 
$Df = df^j \partial_j$ where $df^j$ is a $(0,1)$ tensor and
$\partial_j$ is the pull-back of the coordinate fields on $N$.  
By repeatedly applying the product rule, we have, for all $l$, 
\[\ensuremath{D^{l} [{f}] {} } = \ensuremath{D^{l-1} [{df^j}] {} }\partial_j + 
\sum_{\substack{Z \in \ensuremath{\mathcal{PR}}_{2}(l-1)\\ r_1 \neq l-1}} 
\ensuremath{D^{r_1} [{df^j}] {} } \otimes 
\ensuremath{D^{r_2} [{\partial_j}] {} }
\]
where the derivative on the coordinate vector fields on $N$ is the
pull-back derivative. In any event, the terms in the sum vanish
by applying Lemma~\ref{tl:flat}.
Applying Corollary~\ref{tc:imp}, we find that we have the Euclidean
derivatives of $Df$ when we evaluate the derivative acting on $v^n$ at
$p$. By viewing the coordinates as providing the linear structure, we
have that this is also what appears in the Taylor formula. Thus, $f$
is covariantly normal.
\end{proof}


\section{The exponential map and its derivatives}
\label{s:exp} 

Let $M$ be a cc-manifold. 
Recall that the exponential map at $p \in M$ maps the
tangent space at $p$ into $M$, i.e.\ $\ensuremath{\exp_{p}{}}:\ensuremath{T_p M} \to M$. It
is actually restricted to a neighborhood about the origin such that
its restriction is
a diffeomorphism. Its inverse, \ensuremath{\exp^{-1}_{p}{}}, is a map from an open subset of
$M$ into a flat space.


\begin{theorem} \label{tt:exps}
Let $M$ be a cc-manifold. Then for every $p \in M$, 
we have, for $n >1$, that  $\ensuremath{D_{p}^{n} {[\ensuremath{\exp^{-1}_{p}{}}]} ({v^{n}}) } = 0$ 
and $\ensuremath{D_{0}^{n} {[\ensuremath{\exp_{p}{}}]} ({v^{n}}) } = 0$. For $n=1$, the maps are
essentially the identity.
\end{theorem}

\begin{proof}
By definition, $\ensuremath{\exp^{-1}_{p}{}}:M \to \ensuremath{T_p M}$ is defined such that, when expressed in
normal coordinates, it is the identity map. By Theorem~\ref{t:sim},
we have that  $\ensuremath{D_{p}^{} {[\ensuremath{\exp^{-1}_{p}{(p)}}]} ({v}) } = v$  and, for $n > 1$,  
$\ensuremath{D_{p}^{n} {[\ensuremath{\exp^{-1}_{p}{(p)}}]} ({v^n}) } = 0$. It is immediate that 
$\ensuremath{D_{p}^{} {[\ensuremath{\exp_{p}{(0)}}]} ({v}) } = v$.

To obtain the $n>1$ result for \ensuremath{\exp_{p}{}}, we shall use the chain rule. We
start with the defining relation $\ensuremath{\exp_{p}{(\ensuremath{\exp^{-1}_{p}{(q)}})}} = q$. Thus,  
$\ensuremath{D^{n} [{(\ensuremath{\exp_{p}{\ensuremath{\exp^{-1}_{p}{}}}})}] {} } = \ensuremath{D^{n} [{\mathrm{Id}}] {} }=0$. 
We expand using the chain rule. 
We end up with 
\[
0 = \ensuremath{D_{0}^{n} {[(\ensuremath{\exp_{p}{(\ensuremath{\exp^{-1}_{p}{}})}})]} ({v^n}) } 
=  \sum_{k =1}^n \ensuremath{D^{k} [{\ensuremath{\exp_{p}{}}}] ({\sum_{\ensuremath{Z} \in \ensuremath{\mathcal{\ensuremath{Z}}}(n,k)}
{\mathop{\otimes}_{r \in R_{\ensuremath{Z}}} \ensuremath{D^{|r|} [{\ensuremath{\exp^{-1}_{p}{}}}] ({ 
 v^{|r|}}) }}}) }.
\]
If $ |r| >1$, then  $\ensuremath{D^{|r|} [{\ensuremath{\exp^{-1}_{p}{}}}] ({  v^{|r|}}) } =0$ as
we just established. Just as importantly, if $ |r| =1$, then $\ensuremath{D^{|r|} [{\ensuremath{\exp^{-1}_{p}{}}}] ({
v^{|r|}}) } =v$. The only cr-diagram to survive is the one such that each
row has exactly one column. 
We therefore find that the only term we have
left is $\ensuremath{D^{n} [{(\ensuremath{\exp_{p}{\ensuremath{\exp^{-1}_{p}{}}}})}] ({v^n}) }$ which therefore must
be zero.
\end{proof}

To put it succinctly, the symmetric part of the 
higher derivatives of the exponential map
vanish at the point whose tangent space is being exponentiated. 

Although we do not need the following, it may be useful to remark that
the antisymmetric part of the second derivative of the exponential map
is the torsion. The third derivative decomposes into the symmetric
part, the derivative of the torsion, and the curvature tensor. Both
facts 
can be seen by writing the map in coordinates and looking at the appropriate
symmetry types. We should also mention that if the space is flat, all
the derivatives of the exponential map vanish everywhere. As before,
it is not clear how to prove this statement directly. 

\section{Geodesic submanifolds}

We need one more ingredient before we can give much broader 
sufficient conditions
for covariant normality. In this section, we discuss submanifolds that
 have a connection on them which agree with the ambient manifold's
connection. It
is rare for a submanifold to have this property. For example, the
sphere in \ensuremath{\mathbb{R}^{3}}\ does not have this property. For a Riemannian manifold,
 submanifolds inherit a metric and a connection 
although the inherited connection may be quite different as is the
case  for the sphere as the 
flat metric in \ensuremath{\mathbb{R}^{3}}\ leads to a curved metric on the sphere. 


For the rest of this section, let $M$ be a cc-manifold. 


\begin{defn}
A \emph{geodesic submanifold} $Q$ of a cc-manifold $M$ is a
submanifold which carries a connection that agrees with that of $M$. 
In other words, if $v\in \ensuremath{\Gamma \{ TQ \}}$ and $X \in T_pQ$, then
$\ensuremath{{{\nabla}}_{X}^{Q}{[v]}} = \ensuremath{{{\nabla}}_{X}^{M}{[v]}} \in T_pQ$.
\end{defn}

This notion is
explained in \cite{On83} under the name of totally geodesic submanifold and
in the context of semi-Riemannian manifolds. 
At the heart of the definition is the idea that if $\ensuremath{{{\nabla}}_{x}^{M}{[v]}}
\in T_pQ$ for every $v \in TQ$, then the connection on $Q$ making it
into a geodesic submanifold is just the restriction. In contrast, if
the connection on $M$ does not preserve the tangent space of $Q$ in
the above sense, then one could try to fix this up by using a
projection operator onto the tangent space of $Q$. But, of course,
there is a choice of such projections; a metric does provide such a
choice. 

\begin{theorem}
$Q$ is a geodesic submanifold iff  given any curve $\gamma$ in $Q$, 
parallel transport along $\gamma$ using $M$'s connection is a map
from $T_{\gamma(0)}Q$ into $T_{\gamma(t)}Q$. 
\end{theorem}

\begin{proof}
This is almost immediate. If $Q$ is geodesic, then the defining equations for
parallel transport in $M$ along $\gamma$ are exactly the same as
parallel transport in $Q$ along $\gamma$. Thus, the solutions are the
same. As for the other direction, since the covariant derivative 
 may be defined by
the difference quotient using parallel transport, we see that $Q$ is
geodesic.     
\end{proof}

\begin{theorem}
If $Q$ is a geodesic submanifold, then 
\begin{enumerate}
\item for every $v \in T_pQ$, the geodesic in \ensuremath{\gamma_v}\ in $M$ with
initial tangent vector $v$ is contained in $Q$,
at least for a short while; 
\item $Q$ can be recovered locally by the exponentiation of
neighborhoods in its
tangent spaces. 
\end{enumerate}
\end{theorem}

\begin{proof}
The tangent vector field of a curve in $Q$ is contained in its tangent
space. By looking at what the acceleration of a curve is, it is clear
that both $M$ and $Q$ give the same acceleration. Thus, we define the
geodesics using the connection on $Q$ and then note that
these are geodesics under $M$'s connection. 
The exponentiation statement
follows immediately from the geodesic result. 
\end{proof}


\begin{defn}
Let $p \in M$. 
Then a \emph{flat subspace of \ensuremath{T_p M}} is  a subspace $V_p$ such that  
$\ensuremath{D^{n} [{\ensuremath{\exp_{p}{}}}] {} }$ vanishes on $V_p\ensuremath{^{\otimes n} }$.  
\end{defn}


\begin{defn}
A \emph{flat submanifold} is a geodesic submanifold whose tangent spaces are
flat.  
\end{defn}

We mention a few examples.
All 1-dimensional subspaces are flat. A geodesic is a geodesic
submanifold which is flat. 
 $M$ is flat iff every tangent space to $M$ is flat. This is the remark we made
at the end of the previous section. For vector bundles as manifolds, 
as we shall
explain later, the vertical fiber direct summed with a 1-dimensional
horizontal fiber is a flat subspace. This corresponds to the flat
submanifold of a geodesic on the base space coupled with the fiber
above its points. 


\begin{lemma}
If $Q$ is a geodesic submanifold such that the curvature
and torsion 
of the induced connection vanishes everywhere, then $Q$ is a flat
submanifold.
\end{lemma}

\begin{proof}
This follows by regarding $Q$ as a manifold. Hence any tangent space
of $Q$ is flat. 
\end{proof}

Conceptually, one way to establish that a subspace $V_p$ is flat is to
exponentiate a neighborhood about $0$ and ask whether it is a flat
submanifold. If it is, then $V_p$ is flat. Note that the exponential
image need not even be a geodesic submanifold. How well this captures
the picture of a flat subspace is not known to the author.

We conclude with a reasonable lemma which we shall use in the next
section. 

\begin{lemma} \label{tl:curve}
If $\gamma$ is a curve in a geodesic submanifold $Q$, then all of its
derivatives are contained in the tangent bundle of $Q$. 
\end{lemma}

\begin{proof}
By definition, the \ensuremath{{n}^{th}}\ derivative of $\gamma$ is 
$\nabla_{\dot{\gamma}} \cdots \nabla_{\dot{\gamma}}
\dot{\gamma}$. Since this is the connection evaluated on sections of
$TQ$, we have, by assumption of $Q$ being a geodesic submanifold, that
it is contained in the tangent space of $Q$. 
\end{proof}


\section{Covariant normality} \label{s:maptay}

We may now formulate the positive results of covariant normality.


\begin{lemma} \label{tl:sub}
Let $M$, $N$ be cc-manifolds and $f$ be a smooth map from $M$ into
$N$. Then $f$ is covariantly normal at $p$ if, for each $v \in T_pM$, 
the subspace $ V(f,v) \ensuremath{:=} \text{span}\{ \ensuremath{D_{p}^{n} {[f]} ({v^n}) } \} $ is a
flat subspace of $T_{f(p)}N$. 
\end{lemma}

\begin{proof}
Let $\ensuremath{\tilde{f}}: T_pM \to T_{f(p)}N$ be the following 
map induced on the tangent spaces using the exponential
maps and $f$:
 \[\ensuremath{\tilde{f}}(v) \ensuremath{:=} \ensuremath{\exp^{-1}_{f(p)}{(f(\ensuremath{\exp_{p}{v}}))}}.\] 
Thus, $f(q) = \ensuremath{\exp_{f(p)}{(\ensuremath{\tilde{f}}(\ensuremath{\exp^{-1}_{p}{q}}))}}$. 
Essentially by definition, $\ensuremath{D^{n} [{\ensuremath{\tilde{f}}}] {} }$ 
are the Taylor derivatives. Furthermore,
we also know that $\ensuremath{D_{0}^{n} {[\ensuremath{\exp_{f(p)}{(\ensuremath{\tilde{f}})}}]} ({v^n}) } =
\ensuremath{D_{p}^{n} {[f]} ({v^n}) }$ 
since  
\begin{align*}
\ensuremath{D^{n} [{\ensuremath{\exp_{f(p)}{(\ensuremath{\tilde{f}})}}}] ({v^n}) } =& \ensuremath{D^{n} [{f(\ensuremath{\exp_{p}{v}})}] ({v^n}) } 
\\ =&\sum_{\ensuremath{Z} \in \ensuremath{\mathcal{\ensuremath{Z}}}(n)} 
\ensuremath{D^{|\ensuremath{Z}|_R} [{f}] ({\mathop{\otimes}_{r \in R_{\ensuremath{Z}}} \ensuremath{D^{|r|} [{\ensuremath{\exp_{p}{}}}] ({
v^{|r|}}) }}) } \\
 = & \ensuremath{D^{n} [{f}] ({\{\ensuremath{D^{} [{\ensuremath{\exp_{p}{}}}] ({v}) }\}^{n}}) } = \ensuremath{D^{n} [{f}] ({v^{n}}) }
\end{align*}
where the vanishing of the terms is due to applying
Theorem~\ref{tt:exps}.
We therefore only need to compute the derivatives of $\ensuremath{\exp_{f(p)}{(\ensuremath{\tilde{f}})}}$.
We need to show that
\[
\ensuremath{D^{n} [{\ensuremath{\exp_{f(p)}{(\ensuremath{\tilde{f}})}}}] ({v^n}) }= \ensuremath{D^{n} [{\ensuremath{\tilde{f}}}] ({v^n}) }.
\]
We again use the chain rule. In particular, we have that 
\begin{align*}
\ensuremath{D^{n} [{\ensuremath{\exp_{f(p)}{(\ensuremath{\tilde{f}})}}}] ({v^n}) }  = & 
\sum_{\ensuremath{Z} \in \ensuremath{\mathcal{\ensuremath{Z}}}(n)} 
\ensuremath{D^{|\ensuremath{Z}|_R} [{\ensuremath{\exp_{f(p)}{}}}] ({\mathop{\otimes}_{r \in R_{\ensuremath{Z}}} \ensuremath{D^{|r|} [{\ensuremath{\tilde{f}}}] ({
v^{|r|}}) }}) } \\
= &  \ensuremath{D^{} [{\ensuremath{\exp_{f(p)}{}}}] ({ \ensuremath{D^{n} [{\ensuremath{\tilde{f}}}] ({v^n}) }}) }  \\
&\quad + \sum_{i=2}^{n} \ensuremath{D^{i} [{\ensuremath{\exp_{f(p)}{}}}] ({\sum_{\ensuremath{Z} \in \ensuremath{\mathcal{\ensuremath{Z}}}(n,i)}
{\mathop{\otimes}_{r \in R_{\ensuremath{Z}}} \ensuremath{D^{|r|} [{\ensuremath{\tilde{f}}}] ({ v^{|r|}}) }}}) }
\end{align*}
Since $\ensuremath{D^{} [{\ensuremath{\exp_{f(p)}{}}}] {} } = \ensuremath{\text{Id}}$, the first term is what we want.
Thus, the sum must
vanish in order to obtain the result. 
In general, this will not happen as the input into the derivatives of
the exponential map are not symmetric. 
However, by
the hypothesis, the
derivatives  $\ensuremath{D^{i} [{\ensuremath{\tilde{f}}}] ({v^{i}}) }$ are all elements of the flat
subspace $V(f,v)$. Thus, as the \ensuremath{{n}^{th}}\ derivative of the exponential map
vanishes on $V(f,v)\ensuremath{^{\otimes n} }$, we have that the sum vanishes.   
\end{proof}


\begin{lemma}
Fix $p \in M$. 
Assume that for every geodesic $\gamma$ in $M$ starting at $p$, there
is a {\bf flat}
submanifold $Q_{\gamma}$ 
of $N$ such that $f$ maps the image of $\gamma$ into $Q_{\gamma}$, at
least for some interval about $t=0$.  
Then $f$ is
covariantly normal at $p$.  
\end{lemma}

\begin{proof}
Let $v \in T_pM$ be given and $\ensuremath{\gamma_v}$ be the unique geodesic
corresponding to $v$. Let $V_v$ be the vector space at $p$ tangent to the flat
submanifold $Q_{\ensuremath{\gamma_v}}$ which contains the image of $\ensuremath{\gamma_v}$. Note $V_v$
is a flat subspace implying that 
all we need to establish is that, for each $n$, 
$\ensuremath{D^{n} [{f}] ({v^n}) }$ is an element of $V$. 
Note that $f(\ensuremath{\gamma_v})$ is a curve in a flat submanifold. Thus its
derivatives, of all orders, are elements of $V_v$
i.e.\ $\ensuremath{\frac{{d}^{{n}}
{}}{{d}{{t}^{n}}}}[f(\ensuremath{\gamma_v})] \in V_v$. This follows from Lemma~\ref{tl:curve}.
We also have 
\[
\begin{split}
\ensuremath{\frac{{d}^{{n}}
{f(\ensuremath{\gamma_v})}}{{d}{{t}^{n}}}} &= \sum_{\ensuremath{Z} \in \ensuremath{\mathcal{\ensuremath{Z}}}(n)} 
\ensuremath{D^{ |\ensuremath{Z}|_R } [{f}] ({\mathop{\otimes}_{r \in R_{\ensuremath{Z}}}
\ensuremath{D^{|r|} [{\ensuremath{\gamma_v}}] ({\ensuremath{\frac{\partial {}}{\partial {t}}}^{|r|}}) }}) }\\
  &= \ensuremath{D^{n} [{f}] ({\dot{\ensuremath{\gamma_v}}^n}) }
\end{split}
\]
as the higher derivatives of a geodesic vanish. Since $V_v$ is a flat
subspace and $v$ was arbitrary,
the lemma applies and we have covariant normality of $f$. 
\end{proof}


\begin{theorem} \label{t:gcov}
Assume that for every geodesic $\gamma$ in $M$, there is a flat
submanifold $Q_{\gamma}$ 
of $N$ such that $f$ maps the image of $\gamma$ into $Q_{\gamma}$.  
Then $f$ is
covariantly normal. 
\end{theorem}

\begin{proof}
Apply the lemma at every point $p$ in the manifold.
\end{proof}


\begin{cor}
If $N$ is flat, then all maps $f: M \to N$ are covariantly normal. 
\end{cor}

\begin{proof}
This was actually deduced before and is actually used in the
theorem. Nevertheless, as a means to unification, it is convenient to
note that this follows immediately from the theorem. 
\end{proof}

\begin{cor}
If $N =E$ is a vector bundle over $M$ and $f$ is a section, then $f$
is covariantly normal. 
\end{cor}

\begin{proof}
The proof is delayed until the next section. The main insight is that
the flat submanifold that the geodesic $\gamma$ is mapped into is the
submanifold consisting of the fiber over $\gamma$.
\end{proof}

\begin{cor} \label{tc:geo2}
If $f$ maps geodesics into geodesics, then $f$ is covariantly normal.
\end{cor}

\begin{proof}
A geodesic is a flat submanifold and therefore this is an immediate
application of the theorem.
\end{proof}

\begin{cor} \label{tc:scal}
Let $f: M \to N$. Assume that for every $p\in M$ and $v\in \ensuremath{T_p M}$,  
$f$ satisfies $\ensuremath{D_{p}^{2} {[f]} ({v^2}) } = \alpha_2(v)\ensuremath{D_{p}^{} {[f]} ({v}) }$ for
some $\alpha_2: \ensuremath{T_p M} \to \ensuremath{\mathbb{R}^{}}$.  
Then $f$ is covariantly normal. 
\end{cor}

\begin{proof}
This will follow immediately from Lemma~\ref{tl:sub}, where the subspace
will be 1-dimensional and therefore flat. We shall show that 
$\ensuremath{D^{n} [{f}] ({v^n}) } = \alpha_n(v^{n-1}) \ensuremath{D^{} [{f}] ({v}) }$.


We shall use induction. Fix $p\in M$ and $v_p \in \ensuremath{T_p M}$. 
Extend $v_p$ to a vector
field whose covariant derivative vanishes at $p$; this can be done by
parallel transporting along the geodesics originating at $p$. 
We start with  $\ensuremath{D^{n+1} [{f}] ({v^{n+1}}) } \ensuremath{:=}
\ensuremath{{{\nabla}}_{v}^{}{[\ensuremath{D^{n} [{f}] {} }]}}(v^n)$. 
By  using the product rule, we find
\begin{align*}
\ensuremath{{{\nabla}}_{v}^{}{[\ensuremath{D^{n} [{f}] {} }]}}(v^n) 
& = \ensuremath{{{\nabla}}_{v}^{}{[\ensuremath{D^{n} [{f}] {v^n} }]}} - 
\ensuremath{D^{n} [{f}] {\ensuremath{{{\nabla}}_{v}^{}{[v^n]}}} } \\
& = \ensuremath{{{\nabla}}_{v}^{}{[\alpha_n(v^{n-1}) \ensuremath{D^{} [{f}] ({v}) }]}} - 0 \\
& = \ensuremath{D^{} [{\alpha_n}] ({v^{n-1}, v}) } \ensuremath{D^{} [{f}] ({v}) }  + 0 
+ \alpha_n(v^{n-1}) \ensuremath{D^{2} [{f}] ({v,v}) } +0 \\
&= ( \ensuremath{D^{} [{\alpha_n}] ({v^{n-1}, v}) } + \alpha_n(v^{n-1})
\alpha_2(v))  \ensuremath{D^{} [{f}] ({v}) } \\
& \ensuremath{:=} \alpha_{n+1} (v^n) \ensuremath{D^{} [{f}] ({v}) }.
\end{align*}
This is all we needed to show. Incidentally, the $\alpha$'s are
multilinear functionals. 
\end{proof}

We now show that Corollary~\ref{tc:geo2} and Corollary~\ref{tc:scal} actually
describe the same kind of map. 

\begin{lemma} \cite{EL83}
A map $f: M \to N$ satisfies 
$\ensuremath{D^{2} [{f}] ({[v]^2}) } = \alpha_2(v)\ensuremath{D^{} [{f}] ({v}) }$ if and only if f maps
geodesics of M into geodesics on N. Furthermore,  $\alpha_2$ is
identically zero everywhere if and only if the mapping preserves
parametrizations.  
\end{lemma}

\begin{proof}
The proof is just the chain rule. Let $\ensuremath{\gamma_v}$ be a geodesic on $M$ such that
$\ensuremath{\gamma_v}(0) = p$ and $\ensuremath{{\overset{\cdot}{\gamma}}_{v}({0})} = v$. Then define $\phi$ to be the curve
on $N$ which is the image of $\ensuremath{\gamma_v}$. By the chain rule, we have
\[
\ddot{\phi} = \ensuremath{D^{2} [{f}] ({\dot{\ensuremath{\gamma_v}},\dot{\ensuremath{\gamma_v}}}) } + \ensuremath{D^{} [{f}] ({\ddot{\ensuremath{\gamma_v}}}) }.
\]
Since \ensuremath{\gamma_v}\ is a geodesic, the second term vanishes. The first term
vanishes iff the symmetric part of the second covariant derivative
vanishes. For the full statement, one has to understand that a curve
is a reparametrization of a geodesic iff its acceleration vector is a
multiple of its tangent vector. More precisely, if $\ddot{\phi} (t) =
\alpha(t) \dot{\phi}(t)$, then there is a function $g: \ensuremath{\mathbb{R}^{}} \to
\ensuremath{\mathbb{R}^{}}$ 
such that $\phi(h(s))$ is a geodesic and conversely. The proof follows by
differentiating using the chain rule to arrive at 
\[\ddot{\phi}(h) = \dot{h}^2 \ddot{\phi}(h) + \ddot{h} \dot{\phi}(h).
\]
The only way to have that vanish, without $g$ being a constant, is for
$\ddot{\phi}$ to be a multiple of $\dot{\phi}$ and, furthermore,  $g$
needs to satisfy
$\ddot{h} = - \dot{h}^2 \alpha(t)$, $\dot{h}(0) =1$, $h(0) =0$.
 Assuming that $\alpha$ is
differentiable, we do have local existence of the solution.  That is
all we needed to claim that the symmetric part of the second
 derivative had to be a multiple of the first derivative.
\end{proof}


\section{The failure of covariant normality}

Having established certain situations in which we have covariant
normality, it is rather important to investigate situations in which
we do not have covariant normality which we shall call covariant
abnormality. 
We shall content ourselves with
showing that there are manifolds and maps between those manifolds
which are not covariantly normal. When one looks at what needs to
vanish in order to have covariant normality, one gets the feeling that
a  covariant abnormality should be the generic case.
Unfortunately, it becomes quite complicated to take a brute force
approach  as 
there are many terms that  cancel and/or
are immediately zero. Nevertheless, I would conjecture that given a
manifold map $f : M \to N$, if there exists a tangent 
vector $v \in
T_pM$ such that the 
span of $\{D^n[f] (v)\}$ is not a flat subspace, then $f$ is 
covariantly abnormal at $p$. This is a conjecture with no evidence, but
it seems a natural guess. 

To establish what we want, we shall use normal coordinates. Let $f: M
\to N$ be a map between cc-manifolds. Let $p \in M$ and take a
neighborhood about $p$ small enough so that it 
 can be identified with the tangent space at $p$. Take a basis 
$\{ \partial_i\} $ of the tangent space leading to a coordinatization
 of the neighborhood in $M$.
  We shall use the summation
convention of summing over repeated indices and shall omit summation
symbols. Then the coordinate representation of $f$ is $f \ensuremath{=:} f^i
\partial_i$. Then the first four covariant derivatives of $f$ expressed in
this coordinate system and evaluated on a given vector $v$, are

\begin{align*}
D[f] (v) =& 
D[f^i](v) \partial_i
\\
D^2[f] (v^2) =& 
D^2[f^i](v^2) \partial_i 
+D[f^i](v) D[\partial_i](D[f](v))
\\
\begin{split}
D^3[f] (v^3) =& 
D^3[f^i](v^3) \partial_i 
+2D^2[f^i](v^2) D[\partial_i](D[f](v))\\
&+D[f^i](v) D[\partial_i](D^2[f] (v^2))
+D[f^i](v) D^2[\partial_i](\{D[f](v)\}^2)
\end{split}
\\
\begin{split}
D^4[f] (v^4) =& 
D^4[f^i](v^4) \partial_i 
+3D^3[f^i](v^3) D[\partial_i](D[f](v)) 
\\&
+3D^2[f^i](v^2) D[\partial_i](D^2[f] (v^2))
\\&
+3D^2[f^i](v^2) D^2[\partial_i](D[f](v)\otimes D[f](v))
\\&
+D[f^i](v) D^2[\partial_i](2D^2[f](v^2)\otimes D[f](v) +D[f](v)\otimes
D^2[f](v^2))
\\&
+D[f^i](v) D[\partial_i](D^3[f](v^3))
+D[f^i](v) D^3[\partial_i](\{D[f](v)\}^3) 
\end{split}
\end{align*}

In order to have covariant normality, we need all of the terms to
vanish except the first term, i.e.\ the only term that should survive
is what one would have from Euclidean differentiation in this
coordinate system. 

In order to clean this up a bit, we shall assume that, at $p$, 
 $D[f](v) \ensuremath{=:} \partial_1$ and $D^2[f](v^2) \ensuremath{=:}
 \partial_2$.\footnote{The only loss of generality is assuming that
 $ D^2[f](v^2)$ is not a multiple of $D[f](v)$. If we cannot find a
 point $p$ and a vector $v \in T_pM$ which satisfies this assumption,
 then the map
 is, in fact, covariantly normal.}  
Then the above equations become
\begin{subequations}
\begin{align}
D_p[f] (v) =& \partial_1
\\
D^2_p[f] (v^2) =& \partial_2
\label{te:fcov2}
\\
\begin{split}
D^3_p[f] (v^3) =& 
D^3_p[f^i](v^3) \partial_i 
+2 D_p[\partial_2](\partial_1)\\
&+D_p[\partial_1](\partial_2)
+ D^2_p[\partial_1](\partial_2^2)
\end{split}
\label{te:fcov3}
\\
\begin{split}
D^4_p[f] (v^4) =& 
D^4_p[f^i](v^4) \partial_i 
+3D^3_p[f^i](v^3) D[\partial_i](\partial_1) 
\\&
+3 D_p[\partial_2](\partial_2)
+3 D^2_p[\partial_2](\partial_1^2)
\\&
+ D^2_p[\partial_1](2\partial_2\otimes \partial_1 + \partial_1\otimes
\partial_2)
\\&
+ D_p[\partial_1](D_p^3[f](v^3))
+D^3_p[\partial_1](\{\partial_1\}^3) 
\end{split}
\label{te:fcov4}
\end{align}
\end{subequations} 
 
Please note that since $D_p[\partial_1] (\partial_1) = 0$ at $p$, we
have that \eqref{te:fcov2} implies
\[
D^2_p[f^i] (v^2)= 
\begin{cases} 
0 & \text{if $i \neq 2$} \\ 
1 & \text{if $i = 2$}
\end{cases}  . 
 \]
Furthermore, we have that $D^n_p[w] (w^n) = 0$ for any vector field $w$
which is a sum of the coordinate vector fields with constant
coefficients. This is what normal coordinates accomplish. Taking this
into account in the above expressions, we are left with 
\begin{subequations}
\begin{align}
\begin{split}
D^3_p[f] (v^3) =& 
D^3_p[f^i](v^3) \partial_i 
\\&
+ D_p[\partial_2](\partial_1)
+ D^2_p[\partial_1](\partial_2^2)
\end{split}
\label{te:ffcov3}
\\
\begin{split}
D^4_p[f] (v^4) =& 
D^4_p[f^i](v^4) \partial_i 
+3D^3_p[f^i](v^3) D[\partial_i](\partial_1) 
\\&
+3 D^2_p[\partial_2](\partial_1^2)
+ D^2_p[\partial_1](\partial_2\otimes \partial_1)
\\&
+ D_p[\partial_1](D_p^3[f](v^3)) 
\end{split}
\label{te:ffcov4}
\end{align}
\end{subequations} 
From this point on, we leave generalities. We shall choose $f$ to have
the properties that we wish it to have. Our goal is to establish that
   there exists a map which is not covariantly normal. At some point,
   we shall also make a further assumption on the manifold, but that
   has not come yet. The argument is consideration of certain cases;
   this is a bit distasteful and may largely be the result of the
   author's ignorance as to what constraints are actually placed upon
   the derivatives of these objects. Nevertheless, whatever the facts
   are, we shall be successful in establishing that which we have
   claimed. 


We first show that if torsion does not vanish, then we can choose $f$
such that it is covariantly abnormal. In normal coordinates, we have
that $\tau(x, y) = 2D[x] (y)$ where $x$ and $y$ are sums of the
coordinate vector fields (constant coefficients), 
and the equality is only up to a
sign.\footnote{The ingredients to establish this are:
 $\tau(x,y) \ensuremath{:=} D[x](y) - D[y](x) + [x,y]$; 
  the Lie bracket of coordinate vector fields vanish; and
   normal coordinates imply that the fully symmetric expression for
   covariant derivatives of the coordinate vector fields vanish which
   in turn implies that $D[x] (y) = -D[y](x)$. }
Strangely, the most efficient analysis of the cases involves focussing
on the fourth derivative and assuming that the second covariant derivative
vanishes at $p$ when evaluated at $v$. 
When one makes this assumption, then the third
covariant derivative is the coordinate derivative. The fourth
derivative reduces to 
\[
D^4_p[f] (v^4) = 
D^4_p[f^i](v^4) \partial_i 
+3 D[\partial_3](\partial_1) 
+ D_p[\partial_1](\partial_3) 
\]
where $\partial_3 \ensuremath{:=} D^3[f](v^3)$. 
After applying the vanishing of symmetric expressions, we have 
\[
D^4_p[f] (v^4) = 
D^4_p[f^i](v^4) \partial_i 
+2 D[\partial_3](\partial_1)
\]
Choosing $D_p[f](v) = \partial_1$
and $D^3_p[f](v^3) 
\partial_3$ to be such that the torsion acting on them does not
vanish, we see that the fourth covariant derivative is  not the fourth
coordinate derivative. Hence $f$ is not covariantly normal. Note that
the span of the covariant derivatives is not a flat subspace of the
tangent space. 

If torsion does vanish, then the third covariant derivative will be
the coordinate derivative if and only if 
$D^2_p[\partial_1](\partial_2^2) =0$. I do not know if this vanishes
or not. If it does not, then we are done. If it does vanish, then we must go
to the fourth derivative. For the fourth derivatve to be the
coordinate derivative, we find, after taking into account the
terms we have assumed to vanish,  that we need 
$ D^2_p[\partial_1] (\partial_2 \otimes \partial_1) =0$. Due to the
symmetries already discussed, this is the curvature operator applied
to those vectors. Although such an expression could vanish,
if we assume that $N$ is a curved, two dimensional manifold, then this
cannot vanish. 

As an example, take the sphere and take a curve on the sphere. Then
its acceleration in normal coordinates is the covariant
acceleration. But the higher covariant derivatives of the curve will
not, in general, match  the coordinate derivatives. In particular,
either the third or 
fourth derivative will differ if the acceleration is not zero. Thus,
we have a map which is not covariantly normal. 

As I said in the beginning,  I expect that much stronger statements
can be made, but I do not know how to prove such assertions.


\section{Vector bundles as manifolds} \label{s:vecman}

We now explain why sections of a bundle are covariantly normal. 
 The first item to explain is how a
connection on a vector bundle and the base space give a connection on
the vector bundle as a manifold. Once we have that, then we can view a
section as a map from the base manifold into the bundle manifold and
use linear approximations induced by the normal coordinates. 

We start by recalling the standard fact that on a vector bundle, there
are three equivalent notions of a connection. The first is that of
covariant differentiation and the second is that of parallel
transport. These we have used before. The new one that we need is that
a connection is a choice of horizontal spaces, as we shall now
explain. A reference for this is Addendum 3 of Chapter 8 in \cite{Sp299}. 
Let the vector bundle be $B \ensuremath{:=} (E,\pi, M, V)$  
where $E$ is the total space and it is a manifold;  $M$ is the base
space; $\pi: E \to M$ is the projection map; and $V$ is the fiber,
i.e.\ for every $p$,  $\pi^{-1} (p)$ is isomorphic to $V$.  
For a vector bundle, we also have the (smooth) 
multiplication map: for every $\alpha
\in \ensuremath{\mathbb{R}^{}}$, we have $\tilde{\alpha}: E \to E$ defined by $(p, v)
\mapsto (p, \alpha v)$. For each $\alpha \neq 0$, this mapping is an 
automorphism, i.e.\ a diffeomorphism of $E$ which preserves the bundle
structures. 

Fix a point $e
\in E$ and define $p \ensuremath{:=} \pi(e)$. Then the map $D_e \pi : T_e E \to
T_pM$ is a map between two linear spaces and thus it has a kernel. We
define the  vertical space at $e$, $V_e$, to be the kernel.   
It is natural to try to  decompose
each tangent space into the direct sum $T_e E = V_e \oplus H_e$ where
$H_e$ should be  the horizontal space at $e$. 
We note that any complementary space of $V_e$ is isomorphic 
to $T_p M$ under $D_e \pi$.
The problem is that we have
no canonical way of finding a complementary subspace to the vertical
space. A connection is the assignment of a horizontal space for each
point. 
More
specifically, a distributional connection $H$ in $E$ is a $C^\infty$
distribution such that  $T_eE = H_e \oplus V_e$ 
and that $D_e \tilde{\alpha} (H_e) = H_{\tilde {\alpha}(e)}$ 
for every point $e \in E$
and non-zero $\alpha \in \ensuremath{\mathbb{R}^{}}$. 
As for notation, if $X \in T_e E$, then we define $\tilde{h}_X \in H_e$ and
$\tilde{v}_X \in V_e$ to 
be the horizontal and vertical projections of $X$, respectively; 
thus,  $X = \tilde{h}_X +
\tilde{v}_X$. By using the isomorphism $D_e \pi$, we can view 
$\tilde{h}_X$ as a tangent vector. We shall generally do this and
write $h_X$ for that.  

For vector bundles, as opposed to more general fiber bundles, we are
able to abuse the $v$ term as well. 
Indeed, $V_p \ensuremath{:=} \pi^{-1} (p)$ is a vector space which is also a submanifold
of $E$. At the point $e \ensuremath{:=} (p, v) \in E$, the tangent space of the
fiber is $V_e$. 
But we also know that a vector space is canonically identifiable with
its tangent space. Let $F_e : V_e \to V_p$ be this canonical
isomorphism. We shall use the notation $v_X$ for  $F_e(\tilde{v}_X)$. 

Having setup the basic notations and definitions, we are in a position
to sketch the equivalence of the connections. 
\begin{theorem}
The following notions are equivalent:
\begin{enumerate}
\item
A connection as covariant derivative $(\nabla: T_pM \times
\ensuremath{\Gamma \{ B \}} \to \ensuremath{\Gamma \{ B \}})$ ; 
\item A connection as parallel transport $(P_{\gamma_t}:
V_{\gamma_0} \to V_{\gamma_t} )$; 
\item A connection as a distribution of horizontal spaces $(e \mapsto
H_e)$. 
\end{enumerate}
\end{theorem}

\begin{proof}[Sketch of proof]
We shall only highlight what the source of the relation is; the
details may be found in \cite{Sp299}.

\begin{itemize}

\item[$(1 \Rightarrow 2)$] Given a curve $\gamma$ and a vector $v$ to
be parallel transported, the following ODE, for the function $v_t \in
V_{\gamma_t}$, defines the parallel transport operator: $\dot{v}_t =
0$; roughly written this means $\nabla_{\dot{\gamma}_t} [v(\gamma_t)] =0$.

\item[$(2 \Rightarrow 1)$] Given $v \in \ensuremath{\Gamma \{ B \}}$ and a tangent
vector $y \in T_p M$, let $\gamma$ be a curve through $p$ 
whose tangent vector at $p$ is $y$. Then the covariant derivative is
defined to be $\lim_{t \to 0} \frac{1}{t}
(P_{\gamma_t}^{-1} v(\gamma_t) - v_0)$.

\item[$(1 \Rightarrow 3)$] Let $(p,v) \ensuremath{:=} e \in E$ and
 $\sigma \in \ensuremath{\Gamma \{ B \}}$ such that $\sigma(p) = v$ and $\nabla_y
[\sigma] (p) = 0$ for all $y \in T_pM$. This can be done (at least
locally) by
parallel transporting $v$ along the geodesics emanating
from $p$. Then $H_e$ is the image of $T_pM$ under $D \sigma : T_pM \to T_eE$.  

\item[$(3 \Rightarrow 1)$]
Given a tangent
vector $y \in T_p M$ and a section $\sigma \in \ensuremath{\Gamma \{ B \}}$, 
then $D [\sigma] (y) = \tilde{v}_e + \tilde{h}_e$; we need to have
chosen the horizontal space in order to define
$\tilde{v}_e$.\footnote{Only if $\tilde{h}_e = 0$ would this be
canonical; this, however, never occurs for a section as $p$ is what varies.}  
Thus, the covariant derivative is $v_e$, the canonical isomorphic
image of $\tilde{v}_e$. 

\end{itemize}
\end{proof}


With this information, it is relatively straightforward to setup a
manifold connection on $E$. We shall do this by describing how to parallel
transport tangent vectors.  
Let $\tilde{\gamma}$ be a curve in $E$ going from the point $e_0$ to
$e_1$. Define $\gamma \ensuremath{:=} \pi(\tilde{\gamma})$. This is a curve in
$M$ (although it may only be piecewise smooth, i.e.\ its tangent
vector may vanish). 


Let $X
\in T_{e_0} E$ be the tangent vector we wish to parallel
transport. Let $X = \tilde{h}_X +\tilde{v}_X$, 
the direct sum decomposition of $X$. 
Of course, we now use the isomorphisms to obtain $h_X$ and $v_X$. We
parallel transport those objects along $\gamma$. We 
do this piecewise. We need both connections in this step. Once we have
the transported objects at $e$, we use the isomorphisms to return to a
tangent vector of E. Schematically, we have
\[
X \overset{H_{e_0} \oplus V_{e_0}}{\longmapsto} \tilde{h}_X + \tilde{v}_X 
\overset{D_{e_0} \pi, F_{e_0}}{\longmapsto} (h_X, v_X)
\overset{\gamma}{\mapsto} (g \ensuremath{:=} P_{\gamma}^M h_X, w \ensuremath{:=} P_{\gamma}^B v_x)
{\mapsto} \tilde{g}+
\tilde{w} \in T_{e_1} E  .
\]
One can easily verify linearity, smoothness, invertibility, etc., of
this operation which we define to be the parallel transport operator 
$P_{\gamma}^E$. 

What are the geodesics of $E$? The short answer is that these are the
parallel transport of the bundle's vectors along the geodesics of $M$. 

\begin{lemma}  Let $e = (p,z) \in E$, 
 $X \in T_eE$ and $\gamma$ the geodesic in $M$ with tangent vector
$h_X$.
 Then the
curve defined by $\alpha(t) = \ensuremath{P_{{p},{\gamma(t)}}({tv_X + z})}$ is a
geodesic starting at $e$ whose initial tangent vector is $X$.
\end{lemma}

\begin{proof} 
We shall use the following coordinates. Let $U$ be a 
neighborhood in $M$ centered at $p$ and use normal coordinates. 
Parallel transport  $V_p$ along the geodesics emanating from $p$. This
will locally trivialize the bundle.


With these coordinates, we have $\alpha(t) = (th_X, tv_X +z)$. 
Its time derivative is
$\dot{\alpha}(t) = (h_X, v_X)$. We must now ask whether this vector is
the parallel transport of $(h_X, v_X)$ along $\alpha$. By definition of
the coordinate system, $v_X$ is transported to $v_X$ 
along the geodesic $\gamma(t) = th_X$.  
And the tangent vector of a geodesic is the parallel transport of
itself along the geodesic, so $h_X$ is correct as well. And this is
exactly how we defined parallel transport for the manifold $E$. 
\end{proof}


Thus, we  know the normal neighborhood about a point $e \in
E$ and the exponential map. In fact, we have $\exp{(p,z)}{(\tilde{h},0)}$ is
the parallel transport of $z$ along the geodesic starting at $p$ with
tangent vector $h$. This is the same as $\exp{(p,0)}{(\tilde{h},
\tilde{z})}$. 

We may finally establish covariant normality of sections of $B$ using
the viewpoint that they are maps from $M$ into $E$. Let $\sigma \in
\ensuremath{\Gamma \{ B \}}$. We need to show that it maps geodesics of $M$ into flat
submanifolds of $E$. Then we may apply \ref{t:gcov}. 

\begin{lemma}
Let $Q$ be a geodesic submanifold of $M$. Then $E_Q \ensuremath{:=} \cup_{q
  \in Q} E_q$ is a geodesic submanifold of $E$. 
\end{lemma}

\begin{proof}
First note that $E_Q$ is a submanifold as well as a subbundle of $B$.
Let $\tilde{\alpha}$ be a curve in $E_Q$ and let $\alpha$ be its
projection to $M$ which is actually a curve in $Q$. Parallel transport
of a tangent vector of $Q$ along $\alpha$ stays within the tangent
bundle of $Q$ by assumption of $Q$ being geodesic. Parallel
transporting a vector in $E_p$ to $q$ is definitely an operation which
takes the vector into an element of $E_q$. By looking at the
definition of parallel transport of tangent vectors to $E$, we see
that the parallel transport does preserve  $T E_Q$; thus $E_Q$ is geodesic.   
\end{proof}


\begin{lemma}
Let $Q \ensuremath{:=} \gamma$ be a geodesic. Then $E_Q$ is a flat submanifold. 
\end{lemma}

\begin{proof}
We shall show that parallel transport in $E_Q$ is independent of the
path. First note that parallel transport is independent of the
parametrization of the path. Given a path $\tilde{\alpha}$ in $E_Q$,
Then the projection of $\alpha$ is $\gamma$ parametrized in a
different way; in particular, $\alpha$ can change directions along
$\gamma$. The parallel transport operators multiply together and even
can cancel each other if the path goes forwards and backwards. Thus,
the parallel transport of a tangent vector along $\tilde{\alpha}$ is
just the parallel transport of the appropriate objects along $\gamma$
from the beginning point of $\alpha$ to the terminating point.  
\end{proof}

In contrast, if we assume $Q$ is some higher dimensional flat
submanifold of $M$, we do not have this result unless the curvature of
the bundle connection vanishes on $B_Q$. 

Noting that the image of the geodesic $\gamma$,  under the manifold map
$\sigma \in \ensuremath{\Gamma \{ B \}}$, is contained in the flat submanifold 
$E_{\gamma}$, i.e.\ the set of fibers over $\gamma$, we find that sections of
a bundle are covariantly normal by applying \ref{t:gcov}. 

We make a further remark that the \ensuremath{{n}^{th}} 
Taylor approximation of the section
is the following:
\[
T_p^n[\sigma] (q,\ensuremath{\psi}) = P_{\gamma_{p\to q}}^B (\sum_{i=0}^{n}
D_p^i[\sigma] (\{\ensuremath{\exp^{-1}_{p}{q}} \}^i)
\]
where $P_{\gamma_{p\to q}}^B$ is the parallel transport operator along
the unique geodesic connecting $p$ and $q$ in this small
neighborhood and the covariant derivatives of $\sigma$ are the
bundle's covariant derivatives. This claim follows from considering the
normal coordinates for the manifold $E$; these coordinates correspond
to the basis field described in Appendix \ref{a:difgeom}.

\clearpage

\section{Appendix: a useful frame field}

\subsection{A very useful frame field} \label{a:difgeom} 


In this part, we shall explain how to find sections satisfying
constraints on the symmetric part of their covariant derivative. 
Let $M$ be a manifold with connection and $E$ be a vector bundle over $M$
whose fiber is either isomorphic to \ensuremath{\mathbb{R}^{k}} or $\ensuremath{\mathbb{C}}^k$. Fix $p
\in M$. This appendix relies upon the work of Chapter \ref{c:taylor}
and is used in Chapter \ref{c:geom}


Choose a coordinate neighborhood about $p$ such that for every point
$q$ in the neighborhood, there is a unique
geodesic between $p$ and $q$. Choose a basis $\{e_i\}_{i=1}^{k}$
of the fiber $E_p$ for the rank $k$ bundle $E$  
over $M$. Parallel transport these vectors along the geodesics
emanating from $p$, defining a basis field  also denoted by 
$\{e_i\}_{i=1}^{k}$. 

Let $q$ be any point in the neighborhood. Take
\ensuremath{\gamma_v}\ to be the unique geodesic between $p$ and $q$. Then, by
definition, for each $i$, $\nabla_v [e_i] = 0$ along $\ensuremath{\gamma_v}$. Thus,
Corollary~{\ref{tc:da}} applies and we may conclude that the {\bf symmetric}
 part of  $D^n_p [e_i]
=0$ for all $i$.\footnote{It is true that the diagonal part of a
  tensor determines its symmetric values. One proof is based on the
  following. Define $f(a_1, \ldots,a_k) \ensuremath{:=} A(\{\sum_i a_i e_i \}^n)$.
Notice that we know $f$ for all values of the $a_i$'s. Expand the
right-hand side. Taking the appropriate \ensuremath{\text{n}^{\text{th}}}-order partial derivatives
with respect to the 
$a_i$'s and evaluating the derivatives at $(0,\ldots, 0)$, we can
generate any symmetric expression in the $e_i$'s of length $n$ that we
wish.} 
 We emphasize that this is only true at $p$.

Now define $A := \sum_{i=1}^k f^i e_i$ where the $f^i$ are
scalar functions defined on the neighborhood. Then the symmetric part
of the covariant
derivatives of $A$ are the symmetric covariant derivatives of the $f^i$'s. 
Using Corollary~{\ref{tc:imp}}, we know that the symmetric part of the
higher order covariant derivatives of the $f^i$'s are  the
Euclidean derivatives when expressed in 
 normal coordinates. Thus, it is reasonably
easy to specify the symmetric part of various higher order covariant
derivatives. 

The other parts of the covariant derivatives are part of the geometric
structure of the connection. For example, the second covariant
derivative's antisymmetric part is the curvature tensor, at least if
the connection on $M$ has no torsion. More specifically, we claim
\[
\ensuremath{{{\nabla}}_{x}^{}{[\ensuremath{{{\nabla}}_{y}^{}{[A]}}{}]}} - \ensuremath{{{\nabla}}_{y}^{}{[\ensuremath{{{\nabla}}_{x}^{}{[A]}}{}]}}
- \ensuremath{{{\nabla}}_{[x,y]}^{}{[A]}} 
=
D^2[A] (y\otimes x -x \otimes y) - D[A] (\tau(x, y))
\]
where $\tau$ is the torsion; torsion, acting on a function $f$, is 
\[
\tau(x,y)f = (\ensuremath{{{\nabla}}_{x}^{}{[y]}} - \ensuremath{{{\nabla}}_{y}^{}{[x]}} -[x,y])f 
= D^2[f](x\otimes y -y \otimes x).
\]
These assertions may be verified using the chain rule. A typical fact
to use, for example, is
\[
\ensuremath{{{\nabla}}_{x}^{}{[\ensuremath{{{\nabla}}_{y}^{}{[A]}}{}]}} =  
D^2 [A] (y \otimes x) + D[A](\ensuremath{{{\nabla}}_{x}^{}{[y]}}).
\]
The formulas involving the differentiation of the vector fields are
the usual formulas for the curvature and torsion of a
connection. Writing it in this fashion demonstrates that curvature is
independent of the connection on the manifold. But writing it using
the higher covariant derivatives allows us to immediately see that the
objects are tensors in the tangent vector slots involving $x$ and
$y$. As for the curvature being a tensor in $A$, that is why the
torsion term is necessary. Indeed, multiply $A$ by a function $f$, and
one sees that cross terms cancel when computing the antisymmetric part
leaving the expression in $A$ that we want plus a torsion
term. Subtracting the torsion therefore eliminates the contribution of
$f$. 

Another way to view the subtraction of the torsion
 is to think of torsion as a type of
curvature and that multiplication of a scalar function times a section
is actually the tensor product of the section of $E$ with a section of
the trivial one-dimensional bundle over $M$. Then, as is easily seen,
the curvature of the tensor product connection  is a sum of the two
 curvatures. Thus, we subtract the function's curvature to obtain the
 curvature of the bundle $E$. 

In any event, when we use our special basis, we are only left with the
antisymmetric part  as the symmetric part of
the basis fields vanish at $p$. Thus the second covariant derivatives of these
fields is the curvature tensor applied to field that we are
differentiating, assuming that there is no torsion. 

As an application of this, we can prove that if $E$ is a complex
bundle with a connection and a parallel inner product $\ensuremath{({}, {})}$,
then the curvature tensor is anti-self-adjoint. Let $R$ be the
curvature tensor. Then, to establish our claim, we need to show that 
\[
\ensuremath{({\ensuremath{\psi}}, {R\phi})} = - \ensuremath{({R \ensuremath{\psi}}, {\phi})}
\]
for any sections $\ensuremath{\psi}$ and $\phi$ of $E$. Fix a point $p \in M$. 
Given $\ensuremath{\psi}_p, \phi_p \in
E_p$, define $\ensuremath{\psi}$ and $ \phi$  by parallel transporting the value at
$p$ to all the other points in an appropriately small neighborhood
about $p$. As just explained, $D^2 \ensuremath{\psi} = R\ensuremath{\psi}$ and similarly for
$\phi$ where we use a torsionless connection on $M$ to compute the
second covariant derivative. Since the inner product is parallel, we
have that $\ensuremath{({\ensuremath{\psi}}, {\phi})}_q = \ensuremath{({\ensuremath{\psi}}, {\phi})}_p$ for all $q$ in the
neighborhood. More to the point, we have $D[\ensuremath{({\ensuremath{\psi}}, {\phi})}] =0$ and,
more importantly, 
\[
D^2 [\ensuremath{({\ensuremath{\psi}}, {\phi})}] = 0.
\]
But notice that 
\[
D^2_p[\ensuremath{({\ensuremath{\psi}}, {\phi})}] = \ensuremath{({D_p^2[\ensuremath{\psi}]}, {\phi})} + \ensuremath{({\ensuremath{\psi}}, {D_p^2[\phi]})}
= \ensuremath{({R \ensuremath{\psi}}, {\phi})}_p + \ensuremath{({\ensuremath{\psi}}, {R\phi})}_p;
\]
note that only at $p$ do we have the vanishing of the first derivative
terms in addition to the second derivative terms being the curvature tensor. 
As one can see, we have our result.

\clearpage
\bibliographystyle{abbrv}
\bibliography{references}
\end{document}
